\begin{figure}[tbp]
\centering
\begin{tikzpicture}[x=1cm,y=1cm]
\node[box,text width=2.6cm] (v01) at (0,0) {VMCS01\\L0 运行 L1};
\node[box,text width=2.6cm] (v12) at (4.6,0) {VMCS12\\L1 描述 L2};
\node[greenbox,text width=2.6cm] (v02) at (9.2,0) {VMCS02\\L0 实际运行 L2};
\draw[wire] (v12)--node[above,font=\footnotesize]{状态与控制合成}(v02);
\draw[dwire] (v01.south)--++(0,-0.6)-|node[pos=0.4,below,font=\footnotesize]{同时满足 L0 的控制要求}(v02.south);
\node[note,anchor=west] at (-1.3,-1.45) {逻辑地址映射};
\node[box,text width=1.9cm] (lva) at (0,-2.2) {L2 VA};
\node[box,text width=1.9cm] (lgpa) at (3.4,-2.2) {L2 GPA};
\node[box,text width=1.9cm] (l1pa) at (6.8,-2.2) {L1 GPA};
\node[box,text width=1.9cm] (hpa) at (10.2,-2.2) {host PA};
\draw[wire] (lva)--node[above,font=\footnotesize]{L2 页表}(lgpa);
\draw[wire] (lgpa)--node[above,font=\footnotesize]{EPT12}(l1pa);
\draw[wire] (l1pa)--node[above,font=\footnotesize]{EPT01}(hpa);
\node[note,anchor=west] at (-1.3,-3.1) {实际硬件执行路径};
\node[greenbox,text width=1.9cm] (ava) at (0,-3.9) {L2 VA};
\node[greenbox,text width=1.9cm] (agpa) at (3.4,-3.9) {L2 GPA};
\node[greenbox,text width=1.9cm] (ahpa) at (10.2,-3.9) {host PA};
\draw[wire] (ava)--node[above,font=\footnotesize]{L2 页表}(agpa);
\draw[wire] (agpa)--node[above,font=\footnotesize]{合成 EPT02}(ahpa);
\end{tikzpicture}
\caption[Turtles 的控制与地址映射合成]{Turtles 的控制与地址映射合成。VMCS/EPT 下标表示管理层和被管理层；EPT02 由 L0 按需建立并维护一致性。\cite{source7}}
\label{fig:nested}
\end{figure}
